% ============================================================
% Introduction
% ============================================================
\section{引言 / Introduction}

多模态推理（Multimodal Reasoning）旨在使智能系统能够整合来自文本、图像、视频等多种模态的信息，进行复杂的逻辑推断与决策。近年来，以 GPT-4o、Claude 3 和 Gemini 为代表的大语言模型在视觉问答、图文理解和跨模态生成等任务中取得了显著进展~
cite{alayrac2022flamingo, liu2023llava}。然而，现有模型在处理需要多步推理的复杂场景时，仍面临着模态对齐不稳定、幻觉现象频发以及可解释性不足等挑战~
cite{yin2023survey}。

在本报告中，我们聚焦于以下三个核心研究问题：
\begin{enumerate}
  \item \textbf{模态间语义鸿沟}：如何建立视觉特征与语言概念之间的稳定映射关系？
  \item \textbf{渐进式推理机制}：如何使模型在复杂任务中具备类似人类的逐步推导能力？
  \item \textbf{领域适应性}：如何将通用多模态模型有效迁移至医学影像分析等专业领域？
\end{enumerate}

针对上述问题，我们提出了一种融合视觉-语言表征的渐进式推理框架 PMR。该框架的核心贡献包括：（1）引入了跨模态分层注意力机制，实现细粒度的视觉-语言对齐；（2）设计了基于链式思维（Chain-of-Thought, CoT）的渐进推理模块，显式建模中间推理步骤；（3）构建了面向医学影像的 SIAT-MR 数据集，为领域适应研究提供基准支持。

本报告的结构安排如下：第~\ref{sec:related}~节介绍相关研究工作；第~\ref{sec:method}~节详细阐述 PMR 框架的方法论；第~\ref{sec:experiments}~节汇报实验设置与结果分析；第~\ref{sec:conclusion}~节总结全文并展望未来研究方向。
